\documentclass[12pt, twoside]{article}
%\documentclass[12pt, twoside]{article}
\usepackage[letterpaper, margin=1in, headsep=0.2in]{geometry}
\setlength{\headheight}{0.6in}
%\usepackage[english]{babel}
\usepackage[utf8]{inputenc}
\usepackage{microtype}
\usepackage{amsmath}
\usepackage{amssymb}
%\usepackage{amsfonts}
\usepackage[nomessages]{fp} %\FPeval{\var-name}{2*sin(pi/6)}
\usepackage{siunitx} %units in math. eg 20\milli\meter
\usepackage{yhmath} % for arcs, overparenth command
\usepackage{tikz} %graphics
\usetikzlibrary{quotes, angles, arrows, arrows.meta}
\usepackage{pgfplots}
\usepgfplotslibrary{fillbetween}
\pgfplotsset{compat=1.16}
\usepackage{graphicx} %consider setting \graphicspath{{images/}}
\usepackage{parskip} %no paragraph indent
\usepackage{enumitem}
\usepackage{multicol}
\usepackage{venndiagram}

\usepackage{fancyhdr}
\pagestyle{fancy}
\fancyhf{}
\renewcommand{\headrulewidth}{0pt} % disable the underline of the header
\raggedbottom
\hfuzz=2mm %suppresses overfull box warnings

\usepackage{hyperref}

\fancyhead[LO]{La Scuola d'Italia / Dr.\ Huson / IB Math 12 \\* 2 November 2026}
\fancyhead[RO]{\fbox{\begin{minipage}{6cm}\footnotesize Nickname:\\[0.5cm]\end{minipage}}}
\fancyfoot[C]{\thepage}

\begin{document}
\subsubsection*{3.1 Classwork: The Unit Circle and Trigonometric Graphs (no calculator)}

All angles are in radians. Give exact answers.

\begin{enumerate}[itemsep=0.15cm]

%--- Problem 1: Label the unit circle from memory (16 points); exact values; Pythagorean identity, verify and use with quadrant [14 marks] ---
%--- Source: Original (AA SL 3.4, 3.5, 3.6) ---
\item[1.]

The diagram in the box shows the unit circle with $16$ points marked.

\begin{enumerate}[itemsep=0.1cm]
  \item From memory, write the angle $\theta$ in radians next to each
  point, inside the circle, and its coordinates
  $\left(\cos\theta,\ \sin\theta\right)$ outside the circle.
  \item Write down the exact value of
  $\sin\dfrac{5\pi}{6}$, \ $\cos\dfrac{4\pi}{3}$ \ and \ $\tan\dfrac{7\pi}{4}$.
  \item Verify the identity $\sin^{2}\theta + \cos^{2}\theta = 1$ for
  $\theta = \dfrac{5\pi}{6}$.
  \item Given that $\sin\theta = \dfrac{2}{3}$ and
  $\dfrac{\pi}{2} < \theta < \pi$, find the exact value of $\cos\theta$.
\end{enumerate}

\begin{tikzpicture}
  \draw (-0.6,0) rectangle (14.6,14.4);
  % unit circle, centre (5.0, 9.6), radius 3.2 cm
  \begin{scope}[shift={(5.0,9.6)}]
    \draw[->] (-4.6,0) -- (4.8,0) node[right] {$x$};
    \draw[->] (0,-4.4) -- (0,4.5) node[right] {$y$};
    \draw[thick] (0,0) circle (3.2);
    \foreach \a in {0,30,45,60,90,120,135,150,180,210,225,240,270,300,315,330}
      \fill (\a:3.2) circle (2.2pt);
    \node[below left, font=\footnotesize] at (0,0) {$O$};
  \end{scope}
  \foreach \y in {1,2,3,4} \draw[dotted] (0,\y)--(14.1,\y);
  \foreach \y in {5,6,...,13} \draw[dotted] (10.6,\y)--(14.1,\y);
\end{tikzpicture}

\newpage

%--- Problem 2: Plot sine and cosine from unit-circle values; gradient of OP is tan(theta); line at pi/6 with intercept; line at 2pi/3 through a point [11 marks] ---
%--- Source: Original (AA SL 3.4, 3.5, 2.1) ---
\item[2.]

Use your values from Problem 1. You may use
$\dfrac{\sqrt{3}}{2} \approx 0.87$ and $\dfrac{\sqrt{2}}{2} \approx 0.71$.

\begin{enumerate}[itemsep=0.1cm]
  \item On the grid, plot $y = \sin x$ at $x = 0,\ \dfrac{\pi}{6},\
  \dfrac{\pi}{3},\ \dots,\ 2\pi$. Join the points with a smooth curve and
  label it $y = \sin x$.
  \item On the same grid, plot and label $y = \cos x$ in the same way.
\end{enumerate}

The point $P\left(\cos\theta,\ \sin\theta\right)$ lies on the unit circle,
and the line $OP$ makes an angle $\theta$ with the positive $x$-axis.

\begin{enumerate}[itemsep=0.1cm]
  \item[(c)] Show that the gradient of the line $OP$ is $\tan\theta$.
  \item[(d)] The line $L_{1}$ makes an angle of $\dfrac{\pi}{6}$ with the
  positive $x$-axis and has $y$-intercept $2$. Find the equation of $L_{1}$.
  \item[(e)] The line $L_{2}$ makes an angle of $\dfrac{2\pi}{3}$ with the
  positive $x$-axis and passes through the point $(1,\ 4)$. Find the
  equation of $L_{2}$.
\end{enumerate}

\begin{tikzpicture}
  \draw (-0.6,0) rectangle (14.6,13.5);
  % sketch grid: 1.1 cm per pi/6 across, 2 cm per unit up; origin at (0.6, 10.7)
  \begin{scope}[shift={(0.6,10.7)}]
    \draw[thin, color=gray, xstep=1.1cm, ystep=1.0cm] (0,-2) grid (13.2,2);
    \draw[thick, ->] (-0.3,0) -- (13.7,0) node[right] {$x$};
    \draw[thick, ->] (0,-2.3) -- (0,2.5) node[right] {$y$};
    \foreach \y/\lab in {-2/-1, -1/-0.5, 1/0.5, 2/1}
      \draw (2pt,\y) -- (-2pt,\y) node[left] {$\lab$};
    \foreach \k/\lab in {1/\frac{\pi}{6}, 2/\frac{\pi}{3}, 3/\frac{\pi}{2},
        4/\frac{2\pi}{3}, 5/\frac{5\pi}{6}, 6/\pi, 7/\frac{7\pi}{6},
        8/\frac{4\pi}{3}, 9/\frac{3\pi}{2}, 10/\frac{5\pi}{3},
        11/\frac{11\pi}{6}, 12/2\pi}
      \node[below] at (1.1*\k,-2.1) {$\displaystyle\lab$};
  \end{scope}
  \foreach \y in {1,2,...,7} \draw[dotted] (0,\y)--(14.1,\y);
\end{tikzpicture}

\end{enumerate}
\end{document}
